\section{Evidence synthesis and meta-analysis}
\label{sec:p15-annex-evidence-synthesis}

\begin{annexitem}{annex:p15.synthesis.scale}{thm:p15-synthesis-scale}{Scale transformations and different targets}
\begin{proof}
Linearity gives \(E(b+cX)=b+cE(X)\) and
\(\operatorname{Var}(b+cX)=c^2\operatorname{Var}(X)\).
The common offset cancels in the difference, so its mean and standard
error are multiplied by \(c\). Dividing by an outcome SD transformed by
the same positive \(c\) cancels that factor. In contrast, two populations
with identical mean differences and outcome SDs 1 and 2 have standardized
differences in ratio two to one; no change of units is involved.

Complementation changes \(p_T-p_C\) to its negative. The odds of
a complemented event are the reciprocal original odds, so their ratio is
reciprocated. For risks \(p_T=2/5,p_C=1/5\), the risk ratio is two,
whereas the complementary ratio is \((3/5)/(4/5)=3/4\), not \(1/2\).
Substitute \(p_T=rp_C\) into the risk difference. Solving the odds
equation \(p_T/(1-p_T)=op_C/(1-p_C)\) gives the asserted probability
conversion. The logarithm identity follows from the addition law for
logarithms and the exponential inverse, for positive \(r_i\).

In the non-collapsibility construction, twice the comparator odds give
\(p_T=(2/10)/(9/10+2/10)=2/11\) in the first stratum, and
\(p_T=1/(1/2+1)=2/3\) in the second. Equal stratum probabilities give
\(p_C=(1/10+1/2)/2=3/10\) and
\(p_T=(2/11+2/3)/2=14/33\). Their odds ratio is
\((14/19)/(3/7)=98/57\). Thus even constant conditional odds ratios
do not equal the marginal odds ratio in general.
\end{proof}
\end{annexitem}

\begin{annexitem}{annex:p15.synthesis.linear}{thm:p15-synthesis-linear}{Fixed covariance and optimal pooling}
\begin{proof}
Finite-sum linearity gives the expectation. Expanding the centered square
gives \(\operatorname{Var}(a^{\mathsf T}Y)=
\sum_{i,j}a_ia_j\operatorname{Cov}(Y_i,Y_j)=a^{\mathsf T}\Sigma a\).
When all means equal \(\theta\), normalization of \(a\) gives unbiasedness.
Positive definiteness makes \(\Sigma\) invertible and
\(d=\mathbf1^{\mathsf T}\Sigma^{-1}\mathbf1>0\).
The stated \(a_*\) is normalized and satisfies
\(\Sigma a_*=\mathbf1/d\). Every other normalized \(a\) can be
written \(a_*+b\) with \(\mathbf1^{\mathsf T}b=0\). Consequently
\[
 a^{\mathsf T}\Sigma a=a_*^{\mathsf T}\Sigma a_*
            +b^{\mathsf T}\Sigma b=1/d+b^{\mathsf T}\Sigma b.
\]
The cross term vanishes, and the last term is positive unless \(b=0\).
This proves minimality and uniqueness. A diagonal matrix with entries
\(v_i\) gives the inverse-variance special case.

A linear combination of a multivariate normal vector is normal. For the
independent case this follows directly from the affine normal-sum theorem
\ref{thm:p15-normal-sums}; the covariance representation gives the general
case. Thus \((a_*^{\mathsf T}Y-\theta)/\sqrt{V_*}\) is standard
normal, reducing to \(\sqrt W(\widehat\theta-\theta)\) under independence.
The event that its absolute value does not exceed \(z_{1-\alpha/2}\)
has probability \(1-\alpha\), proving coverage. The variance expansion
used only second moments. For example, a single symmetric observation
\(Y=\theta\pm1\) with equal probabilities has variance one but does not
have a standard-normal pivot. Normal coverage cannot be inferred from the
variance calculation alone.
\end{proof}
\end{annexitem}

\begin{annexitem}{annex:p15.synthesis.estimated-weights}{thm:p15-synthesis-estimated-weights}{Data-dependent weights and their bias}
\begin{proof}
Expand each \(E(A_iY_i)=EA_i\,EY_i+\operatorname{Cov}(A_i,Y_i)\).
The normalized weights give \(\sum_iEA_i=1\); summing proves the identity.
For independent \(Y_1,Y_2\sim N(\theta,1)\), the difference has
\(N(0,2)\) law and
\[
 E|Y_1-Y_2|=2\int_0^\infty {x\over2\sqrt\pi}e^{-x^2/4}\,dx
                  ={2\over\sqrt\pi}.
\]
Since \(\max(Y_1,Y_2)=(Y_1+Y_2+|Y_1-Y_2|)/2\), its mean is
\(\theta+1/\sqrt\pi\). For any \(0<\varepsilon<1/2\), give the
larger observation weight \(1-\varepsilon\) and the smaller weight
\(\varepsilon\). Both weights are positive and the estimate equals
\(\varepsilon(Y_1+Y_2)+(1-2\varepsilon)\max(Y_1,Y_2)\).
Its mean is \(\theta+(1-2\varepsilon)/\sqrt\pi\). Ties have probability
zero and may be resolved arbitrarily. This proves the positive-weight
counterexample and its limit.

For a chosen fixed weight vector the actual variance remains
\(a^{\mathsf T}\Sigma a\). Substituting a different covariance matrix
into that formula does not change the data law. The duplicated/shared-control
constructions in \ref{annex:p15.synthesis.dependence} give explicit failures
of a diagonal replacement. This argument does not assert that every estimated
weight procedure is biased; it establishes that fixed-weight exactness does
not prove a data-dependent procedure's properties.
\end{proof}
\end{annexitem}

\begin{annexitem}{annex:p15.synthesis.score}{thm:p15-synthesis-score}{Conditional score, information and one Newton step}
\begin{proof}
For fixed arm sizes, let \(C_i\) be the comparator event-count variable.
The independent binomial probabilities of
\(A_i=x,C_i=m_i-x\) are proportional, after conditioning on \(m_i\),
to the two binomial coefficients times
\(\{p_{Ti}(1-p_{Ci})/[p_{Ci}(1-p_{Ti})]\}^{x}\).
All remaining factors depend on the margins but not on \(x\). Thus the
stated odds model has finite normalizing sum \(Z_i(\delta)>0\).
Its log likelihood is \(\delta\sum_i a_i-\sum_i\log Z_i(\delta)\),
up to a constant. Finite-sum differentiation gives
\(Z_i'/Z_i=E_\delta A_i\) and
\((\log Z_i)''=E_\delta A_i^2-(E_\delta A_i)^2\).
These are the asserted score and information.

At zero, every selection of \(m_i\) event positions among \(n_i\)
positions is equally likely. Let \(J_j\) indicate that treatment position
\(j\) is selected. Counting selections gives
\(EJ_j=m_i/n_i\) and, for different positions,
\(EJ_jJ_l=m_i(m_i-1)/(n_i(n_i-1))\).
Since \(A_i\) is the sum over \(n_{Ti}\) such positions, these identities
give \(E_i=n_{Ti}m_i/n_i\) and
\[
 \operatorname{Var}A_i
 =n_{Ti}{m_i\over n_i}\left(1-{m_i\over n_i}\right)
  +n_{Ti}(n_{Ti}-1)
       \left\{{m_i(m_i-1)\over n_i(n_i-1)}-{m_i^2\over n_i^2}\right\}
 =V_i.
\]
Independence between studies makes the null score centered with variance
\(\sum_i V_i\). Newton's update for solving \(U(\delta)=0\) from
zero is \(0-U(0)/U'(0)=U(0)/I(0)\), exactly the displayed one-step
calculation. An approximately quadratic log likelihood motivates its
approximate SE; no quadratic-exactness or normal score law was proved.
If \(m_i=0\), the conditional support contains only \(A_i=0\);
if \(m_i=n_i\), it contains only \(A_i=n_{Ti}\).
Both have zero variance and score. In particular no-events tables contribute
no conditional relative-effect information, while the observed absence of
events remains information about event frequency under other models.
\end{proof}
\end{annexitem}

\begin{annexitem}{annex:p15.synthesis.table-delta}{thm:p15-synthesis-table-delta}{Consistency and first-order table-ratio uncertainty}
\begin{proof}
For each arm, the sample event proportion converges in probability to its
fixed event probability: its variance is \(p(1-p)/n_{gi}\to0\),
so Chebyshev's inequality suffices. Normalize all counts and arm sizes by
\(n\). The normalized Mantel--Haenszel numerator and denominator therefore
converge to positive constants. For the odds ratio these are sums of
\(\lambda_{Ti}\lambda_{Ci}p_{Ti}(1-p_{Ci})/(\lambda_{Ti}+\lambda_{Ci})\)
and
\(\lambda_{Ti}\lambda_{Ci}(1-p_{Ti})p_{Ci}/(\lambda_{Ti}+\lambda_{Ci})\).
When each odds ratio is the same \(o\), every numerator summand is
\(o\) times its denominator summand. Their limiting ratio is \(o\).
For a common risk ratio, the two limits instead contain \(p_{Ti}\)
and \(p_{Ci}\); their ratio is that common risk ratio by the same argument.
At the expected arm counts, the relevant ratio equals the common effect
for every finite set of arm sizes. Consequently slow convergence of the
deterministic arm fractions causes no additional centering error below.

The independent Bernoulli central limit theorem
\ref{thm:p15-clt-iid}, applied to each of the finitely many arms, gives
joint independent normal limits for their centered proportions at the
\(\sqrt n\) scale. Joint convergence follows because their joint
characteristic functions factor into the marginal characteristic functions.
The limiting variances are \(p_{gi}(1-p_{gi})/\lambda_{gi}\).
The log ratio is a continuously differentiable function of those proportions
and arm fractions in a neighborhood of their limiting values: its two
totals are positive there. Its first-order Taylor expansion has a remainder
\(o_{\mathbb P}(n^{-1/2})\). To see this directly, the maximum change
of its gradient over a shrinking neighborhood tends to zero by continuity;
multiply that change by the tight \(O_{\mathbb P}(n^{-1/2})\)
proportion fluctuations. Outside the neighborhood the probability tends to
zero by consistency. The limiting linear combination is centered normal
with variance obtained by the displayed independent-arm covariance expansion.
It has positive variance, since the event probabilities are interior and
the risk/odds comparison has nonzero derivatives.

Differentiating \(\log R-\log S\) with respect to the two observed
event counts gives the four displayed derivative formulas. The binomial
count variances are \(n_{gi}p_{gi}(1-p_{gi})\). Their first-order
covariance expansion yields \(\widehat V_f\) after substitution of
consistent sample proportions and gradients. The resulting variance ratio
converges to one; \ref{thm:p15-clt-perturbation} therefore gives the
standard-normal standardized error. Events on which observed totals vanish
have probability tending to zero under these interior-probability assumptions;
one may assign any fixed finite statistic on those events for the limit
argument. These conclusions do not cover sparse-event sequences whose
probabilities tend to zero, an increasing number of tiny studies, dependent
arms, or exact finite-sample interval coverage.
\end{proof}
\end{annexitem}

\begin{annexitem}{annex:p15.synthesis.change}{thm:p15-synthesis-change}{Repeated measures and a conditional scale bridge}
\begin{proof}
Expand the centered square of \(X_1-X_0\) to obtain its two variances
minus twice their covariance. The difference of group mean changes is
\((EX_{1T}-EX_{1C})-(EX_{0T}-EX_{0C})\); equal population baseline
means remove its second term. If each final SD and baseline SD is one,
their covariance is \(7/8\), and groups have the same sampling sizes,
the change SD is \(\sqrt{2-2(7/8)}=1/2\),
whereas the final SD is one. A mean difference of one therefore standardizes
to two or one respectively. Equality of the MD target does not identify
these standardized targets or their variances.

For threshold \(c\), the exceedance probability in the stated logistic
model is \(1-F_g(c)\). Its odds are
\((1-F_g(c))/F_g(c)=\exp((m_g-c)/s)\).
Taking the ratio between groups cancels the threshold and gives
\(\log\mathrm{OR}=(m_T-m_C)/s\).
Dividing the mean difference by the SD specified in the conversion premise,
\(s\pi/\sqrt3\), gives \((\sqrt3/\pi)\log\mathrm{OR}\).
This is a conditional derivation under the declared latent distribution and
SD premise. It proves no distribution-free bridge between continuous and
dichotomized measurements. A deterministic scalar conversion of a statistic
multiplies its standard deviation by the absolute scalar, proving the stated
SE conversion within that model.
\end{proof}
\end{annexitem}

\begin{annexitem}{annex:p15.synthesis.random-known}{thm:p15-synthesis-random-known}{Mean uncertainty and future-effect uncertainty}
\begin{proof}
Independent normal sums give \(Y_i\sim N(\mu,v_i+\tau^2)\).
The linear pooling result applied with these variances gives normal mean
\(M(\tau^2)\) of variance \(1/W(\tau^2)\).
The new \(\Theta_*\) is independent of the original data, so
\(\Theta_*-M\sim N(0,\tau^2+1/W)\). Standardizing proves the stated
prediction probability, including \(\tau^2=0\). If a future independent
normal error with variance \(v_*\) is added, the same normal-sum computation adds
\(v_*\) to the variance. An interval of radius \(z/\sqrt W\)
uses the variance of estimating \(\mu\) alone, and excludes the extra
between-study variation when \(\tau^2>0\). Independence and the same
future-effect model are essential to this probability statement.
With all \(v_i=v\), the total weight is \(k/(v+\tau^2)\).
Substitution gives both displayed equal-variance formulas and their limits.
\end{proof}
\end{annexitem}

\begin{annexitem}{annex:p15.synthesis.q}{thm:p15-synthesis-q}{Dispersion expectation and Q-profile coverage}
\begin{proof}
The identity
\(Q(0)=\sum_iw_iY_i^2-(\sum_iw_iY_i)^2/W\) follows by expanding
the weighted residuals. Center the independent observations at their common
mean. With variances \(v_i+\tau^2\), the mean terms cancel and
\[
 E Q(0)=\sum_iw_i(v_i+\tau^2)
       -{\sum_iw_i^2(v_i+\tau^2)\over W}
       =k-1+\tau^2\left(W-{\sum_iw_i^2\over W}\right).
\]
The constant \(C=2\sum_{i<j}w_iw_j/W\) is strictly positive for
\(k\geq2\). Solving the moment equation gives the untruncated estimate;
restriction to nonnegative values gives DL as defined. At \(\tau^2=0\),
the untruncated expression has mean zero, while its positive part generally
has positive expectation, so unbiasedness cannot be inherited after truncation.

For \(u\) equal to the true \(\tau^2\), put
\(Z_i=(Y_i-\mu)\sqrt{w_i(u)}\) and
\(d_i=\sqrt{w_i(u)/W(u)}\). The \(Z_i\) are independent standard
normals, \(d^{\mathsf T}d=1\), and
\(Q(u)=Z^{\mathsf T}(I-dd^{\mathsf T})Z\).
Complete \(d\) to an orthonormal basis. Under its orthogonal change of
coordinates the joint normal density is unchanged: the squared norm is
unchanged and the absolute Jacobian is one. Thus the resulting coordinates
are independent standard normals. The quadratic form is the sum of squares
of the \(k-1\) coordinates perpendicular to \(d\), hence has
\(\chi^2_{k-1}\) law by the normal/gamma construction of
\ref{thm:p15-student-orthogonal}.

For a general specified \(u\), the weighted mean uniquely minimizes
\(\sum_iw_i(u)(Y_i-m)^2\) over \(m\). Differentiating its value
at that minimum, the mean-derivative term vanishes because
\(\sum_iw_i(Y_i-M)=0\). Therefore
\(Q'(u)=-\sum_iw_i(u)^2(Y_i-M(u))^2\leq0\).
The exact confidence set contains the true variance precisely when its
chi-squared variable lies between the two indicated quantiles, an event of
probability \(1-\alpha\). For equal observed \(Y_i\), every \(Q(u)\)
is zero and the lower positive quantile makes the set empty; analogous
low-dispersion samples can give an empty set as well. This is compatible with
the stated coverage property. If variances are estimated, the independent
standard-normal whitening used above has not been established, so the exact
chi-squared law and consequent exact coverage do not follow.
\end{proof}
\end{annexitem}

\begin{annexitem}{annex:p15.synthesis.reml}{thm:p15-synthesis-reml}{Restricted likelihood and existence of a maximizer}
\begin{proof}
The normal product likelihood is proportional to
\(\prod_i(v_i+u)^{-1/2}
\exp[-\sum_iw_i(u)(Y_i-m)^2/2]\).
Complete the weighted square:
\(\sum_iw_i(Y_i-m)^2=Q(u)+W(u)(m-M(u))^2\).
Its integral over \(m\) with respect to Lebesgue measure is proportional to
\(\prod_i(v_i+u)^{-1/2}W(u)^{-1/2}e^{-Q(u)/2}\).
This integral represents the likelihood of contrasts that eliminate the
mean, up to a factor independent of \(u\): take an orthonormal coordinate
in direction \(k^{-1/2}\mathbf1\) and \(k-1\) perpendicular
coordinates. Integrating the original density over the first coordinate
is its contrast density; substituting the coordinate by a simultaneous
translation of every \(Y_i\) changes the integration variable by
the constant \(\sqrt k\). Consequently the integrated expression differs
only by that constant. This argument is a restricted-likelihood construction,
not adoption of an improper Bayesian prior.

Taking logs gives \(\ell_R\). Since \(w_i'=-w_i^2\),
\((\log W)'=-\sum_iw_i^2/W\), and the preceding result gives
\(Q'=-\sum_iw_i^2(Y_i-M)^2\), differentiation proves the score formula.
Replacing \(m\) by its maximizing value without integration would omit
\(-\log W/2\), so that profile likelihood is different.

The function is continuous on \([0,\infty)\), because all \(v_i\)
are positive. As \(u\to\infty\),
\(\sum_i\log(v_i+u)=k\log u+o(1)\),
\(\log W(u)=\log k-\log u+o(1)\), and \(Q(u)\to0\).
Thus \(\ell_R(u)=-(k-1)\log u/2+O(1)\to-\infty\).
Its maximum is therefore attained on some compact interval by continuity.
This establishes existence, not uniqueness or an unbiased estimate of \(u\).
\end{proof}
\end{annexitem}

\begin{annexitem}{annex:p15.synthesis.inconsistency}{thm:p15-synthesis-inconsistency}{The two inconsistency calculations}
\begin{proof}
When \(Q(0)\leq k-1\), the DL estimate and both summaries are zero.
When \(Q(0)>k-1\), substitute
\(\widehat\tau^2=(Q(0)-(k-1))/C\) and
\(\widetilde v=(k-1)/C\). Their ratio is
\((Q(0)-(k-1))/Q(0)\), giving the identity after conversion to percent.
With fixed positive \(\widehat\tau^2\), dividing numerator and
denominator of its ratio by that quantity gives
\(1/(1+\widetilde v/\widehat\tau^2)\to1\).
This limit changes relative precision but leaves the absolute variance
unchanged, establishing the interpretive boundary.
\end{proof}
\end{annexitem}

\begin{annexitem}{annex:p15.synthesis.fixed-shape-t}{thm:p15-synthesis-fixed-shape-t}{The exact weighted Student model}
\begin{proof}
Set \(Z_i=(Y_i-\mu)\sqrt{w_i/c}\) and
\(d_i=\sqrt{w_i/W}\). They are independent standard normals and \(d\)
is a unit vector. The coordinate \(d^{\mathsf T}Z\) equals
\(\sqrt{W/c}(M-\mu)\), while the perpendicular squared norm is
\(\sum_iw_i(Y_i-M)^2/c=(k-1)q_H/c\).
The orthogonal density argument in \ref{annex:p15.synthesis.q} makes
the first coordinate independent of the remaining ones. The residual sum
is \(\chi^2_{k-1}\), and their standardized ratio is Student by
\ref{thm:p15-student-law}. Its two-sided quantile event has probability
\(1-\alpha\), as claimed. Data-dependent weights do not produce the
fixed orthogonal transformation of independent standard normals used in this
proof. An approximate HKSJ interpretation needs separate assessment.
\end{proof}
\end{annexitem}

\begin{annexitem}{annex:p15.synthesis.bayes}{thm:p15-synthesis-bayes}{Posterior normalization and future effects}
\begin{proof}
For each fixed observed vector, its normal likelihood is strictly positive
everywhere on \(\mathbb R\times[0,\infty)\). It is bounded above
by \(\prod_i(2\pi v_i)^{-1/2}\), because the exponential factors do
not exceed one and \(v_i+u\geq v_i\). Integrating against the prior
probability measure gives a positive finite normalizer. Dividing by it
defines the posterior, with the stated likelihood kernel. The same bound
and the prior moment assumptions imply finite posterior second moment for
\(\mu\) and finite posterior first moment for \(u\).

Conditional on \((\mu,u)\) and the observations, the independent new
effect has mean \(\mu\), variance \(u\), and second moment
\(u+\mu^2\). Posterior integration gives
\(E(\Theta_*\mid Y)=E(\mu\mid Y)\) and
\(E(\Theta_*^2\mid Y)=E(u\mid Y)+E(\mu^2\mid Y)\).
Subtract the squared posterior mean to obtain the variance identity.
These calculations apply to a prior with atoms as well as a prior density.
If the same data are used twice as if independent, its likelihood factor is
squared rather than counted once. For example, a normal datum with variance
\(v\) and a normal prior of variance \(b\) gives posterior precision
\(1/b+1/v\); duplicating its likelihood gives \(1/b+2/v\), a different
posterior from the same evidence. A model of genuine independent replication
would require another independently generated observation.
\end{proof}
\end{annexitem}

\begin{annexitem}{annex:p15.synthesis.dependence}{thm:p15-synthesis-dependence}{Shared observations and covariance}
\begin{proof}
For exact duplicate reports \(Y_1=Y_2=Y\), their average is literally
\(Y\), so its variance is \(v\). The covariance matrix is
\(v\left(\begin{smallmatrix}1&1\\1&1\end{smallmatrix}\right)\),
not the independent diagonal matrix; it is singular because the two records
contain the same observation.
For independent arm means \(\overline T_1,\overline T_2,\overline C\),
expand
\(\operatorname{Cov}(\overline T_1-\overline C,
\overline T_2-\overline C)\). The cross-arm terms vanish by independence
and the remaining term is \(\operatorname{Var}\overline C\).
Finally \(\operatorname{Cov}(Y_1,Y_2)=rv\) gives
\(\operatorname{Var}((Y_1+Y_2)/2)=(v+v+2rv)/4\).
The stipulated correlation range is exactly the positive-semidefinite
range for that two-by-two equal-variance matrix, including its singular
endpoints. Separate marginal variances do not determine \(r\).
\end{proof}
\end{annexitem}

\begin{annexitem}{annex:p15.synthesis.subgroups}{thm:p15-synthesis-subgroups}{Within and between groups}
\begin{proof}
In each group write \(Y_i-M=(Y_i-M_g)+(M_g-M)\). The cross term
in the weighted square sum vanishes because
\(\sum_{i\in g}w_i(Y_i-M_g)=0\). Summing the remaining two terms
proves the decomposition.
Under independent common-mean normals, each \(M_g\) has variance
\(1/W_g\), and different group means are independent because they use
disjoint independent studies. Apply the projection argument for dispersion
to these \(G\) independent normals with weights \(W_g\); the
between-group term has \(\chi^2_{G-1}\) law.
Expanding the square of the difference of two general centered summaries
gives variance \(V_1+V_2-2c_{12}\).
For the displayed numbers the two individual standardized means are 2.5
and 1.8; their independent contrast statistic is
\((1.0-0.9)/\sqrt{0.4^2+0.5^2}\approx0.15617\).
This explicit failure separates the within-group significance questions
from the question about a difference between groups.
\end{proof}
\end{annexitem}

\begin{annexitem}{annex:p15.synthesis.missing-bounds}{thm:p15-synthesis-missing-bounds}{What missing binary outcomes permit}
\begin{proof}
Among the \(m\) missing binary outcomes the event count is an integer
\(x\in[0,m]\). The full-sample proportion is \((e+x)/N\),
bounded by the stated endpoints. Every integer count in that range is
realized by assigning that many missing outcomes to the event.
For the smallest difference choose no missing treatment events and all
missing comparator events; for the largest choose all missing treatment
events and no missing comparator events. These choices are separately
feasible and attain the displayed limits. For example, if each group has
100 allocated participants, with treatment \(e_T=20,m_T=20\) and
comparator \(e_C=25,m_C=10\), the difference lies in
\([-.15,.15]\). Complete-case proportions \(20/80\) and \(25/90\)
give approximately \(-.02778\), which does not resolve the sign of the
full-sample difference. No sampling-coverage statement was used.
\end{proof}
\end{annexitem}

\begin{annexitem}{annex:p15.synthesis.sensitivity}{thm:p15-synthesis-sensitivity}{Deletion, weights and bias bounds}
\begin{proof}
The deleted numerator is \(WM-w_jY_j\) and denominator \(W-w_j\).
Subtracting \(M\) gives the stated identity. Differentiating the original
ratio with the other weights fixed gives \((Y_j-M)/W\).
For shifts \(b_i\), the estimate's shift is \(\sum_i(w_i/W)b_i\).
Positive normalized weights and the triangle inequality bound its absolute
value by \(\sum_i(w_i/W)B_i\). Choosing every \(b_i=B_i\), or
every \(b_i=-B_i\), attains the two extremes. These are fixed-weight
identities. A heterogeneity estimate obtained from the retained studies
depends on which studies remain, so it generally changes the other weights
and requires a fresh calculation.
\end{proof}
\end{annexitem}

\begin{annexitem}{annex:p15.synthesis.example}{thm:p15-synthesis-example}{Complete arithmetic and failure examples}
\begin{proof}
For \((-2,-1,2)\) and weights \((1,1,1/4)\),
\(\sum w_iY_i=-5/2\), \(\sum w_iY_i^2=6\), and \(W=9/4\).
Thus \(M=-10/9\) and \(Q=6-(25/4)/(9/4)=29/9\).
Also \(\sum w_i^2=33/16\), so
\(C=9/4-(33/16)/(9/4)=4/3\) and
\(\widehat\tau_{\rm DL}^2=(29/9-2)/(4/3)=11/12\).
The common-effect SE is \(2/3\), giving the quoted interval using
\(z_{.975}\approx1.959964\).

At the DL estimate, the two first weights are \(12/23\) and the third
is \(12/59\). Their sum is \(1692/1357\), their weighted numerator
is \(-1572/1357\), and the mean is \(-131/141\).
The weighted residual sum is \(2536/1081\). Hence
\[
 q_H=1268/1081,\quad s_H^2=18703/19881,
 \quad s_M=\sqrt{1357/1692}\approx0.895550,
 \quad s_H\approx0.969921.
\]
The Wald interval is \([-2.68432,0.82617]\). Using
\(t_{2,.975}\approx4.302653\), the HKSJ interval is
\([-5.10231,3.24416]\). The HKSJ-based prediction calculation with
the same t multiplier and \(s_H\) is
\([-6.79303,4.93488]\). The Wald-based prediction calculation with
the normal multiplier and \(s_M\) is \([-3.49856,1.64040]\).
These last three numerical intervals are the stated approximate procedures,
not exact coverage results for these estimated weights. Deleting study 3
with the fixed common-effect weights gives \(-3/2\), with SE
\(1/\sqrt2\); the shift is \(-7/18\), also obtained from the deletion
identity. Re-estimating DL on the two remaining studies gives \(Q=1/2\)
and \(\widehat\tau^2=0\), illustrating why a full random-effects deletion
analysis must recompute its heterogeneity.

The duplicate-record variance computation in
\ref{annex:p15.synthesis.dependence} proves the standard-error error factor
\(\sqrt2\). The subgroup contrast computation in
\ref{annex:p15.synthesis.subgroups} proves the significance-label failure.
For twenty independent null rejection events of probability .05, independence
gives probability \(.95^{20}\) of none. Its complement is
\(1-.95^{20}\approx.641514\). Without independence, subadditivity
gives at most \(20(.05)=1\), not that exact value.

For the aggregation failure, suppose each of two studies consists equally
of a low-age and a high-age group. Both a within-study effect function
\(\theta(x)=x\), \(x\in\{0,1\}\), and the constant function
\(\theta(x)=1/2\) yield study average \(1/2\) and average covariate
\(1/2\). The study-level summaries are identical but the first function
has effect modification and the second does not. Those summaries therefore
cannot identify the unreported within-study modification. Finally, averaging
unrelated measurement targets supplies only their numerical arithmetic;
the scale and non-collapsibility counterexamples above show why an identified
common target must precede that arithmetic.
\end{proof}
\end{annexitem}
